\section{Related Work}\label{sec:related-work}

\subsection{Retrieval-Augmented Generation}\label{sec:rw-rag}

Lewis et al.~\cite{lewis2020rag} introduced RAG as a method for grounding
language model outputs in retrieved passages, establishing the retrieve-then-read
paradigm. Concurrent work on REALM~\cite{guu2020realm} demonstrated that
retrieval augmentation could be integrated into pre-training. Dense Passage
Retrieval (DPR)~\cite{karpukhin2020dpr} established embedding-based retrieval
as the dominant approach, and Fusion-in-Decoder~\cite{izacard2021leveraging}
showed that conditioning generation on multiple retrieved passages improves
multi-hop performance.

More recent work has explored self-reflective
retrieval~\cite{asai2023selfrag}, which trains models to decide when to
retrieve and to critique retrieved content. Gao et al.~\cite{gao2024ragsurvey}
survey the current landscape, cataloguing trade-offs between retrieval
granularity, index size, and generation quality.

A consistent theme across RAG work is that token budget is treated as an
engineering constraint rather than an evaluation axis. The BEIR
benchmark~\cite{thakur2021beir} evaluates zero-shot retrieval across 18
heterogeneous IR tasks, and RAGAS~\cite{es2023ragas} measures faithfulness and
relevance of RAG pipelines---but neither measures per-query token consumption.
This gap motivates our Reasoning Density Score (Section~\ref{sec:metrics}).

\subsection{Multi-Hop Question Answering}\label{sec:rw-multihop}

The SQuAD benchmark~\cite{rajpurkar2016squad} established token-level F1 as
the standard for extractive QA, which we adopt for its well-understood
properties. HotpotQA~\cite{yang2018hotpotqa} introduced multi-hop reasoning
as a distinct challenge: answering a question requires traversing multiple
supporting facts. MuSiQue~\cite{trivedi2022musique} further isolated
compositional multi-hop reasoning from shortcut exploitation.

These benchmarks evaluate whether a system can \textit{find} multi-hop
answers in unstructured text. Our T3 query type tests the related but
distinct problem of \textit{traversing} a structured dependency graph---a
task where the ground truth is the explicit path, not an extracted span.
RAG must infer graph paths from prose; CKG traverses them directly.

\subsection{Graph-Based Retrieval}\label{sec:rw-graphrag}

Edge et al.~\cite{edge2024graphrag} introduced GraphRAG, which applies
community detection over LLM-extracted entity graphs to answer both local
(entity-level) and global (corpus-level) queries. LightRAG~\cite{guo2024lightrag}
reduces GraphRAG's computational footprint with a dual-level retrieval
strategy. HippoRAG~\cite{gutierrez2024hipporag} draws on hippocampal memory
models to build associative knowledge networks from text.

Graph retrieval applied specifically to structured domain representations is
explored in G-Retriever~\cite{he2024gretriever}, which encodes text-attributed
graphs for QA, and Think-on-Graph~\cite{sun2024thinkongraph}, which performs
beam search over knowledge graphs to reason step-by-step. Peng et
al.~\cite{peng2024graphrag2} survey the rapidly growing Graph-RAG literature.

A key distinction for this benchmark: GraphRAG, LightRAG, and HippoRAG all
perform \textit{dynamic extraction}---they derive graph structure from
unstructured text at index time. CKG uses \textit{pre-structured domain
knowledge} in which the graph was authored by a domain expert. When expert
structure is available, the dynamic extraction step is wasted computation;
our benchmark quantifies this waste.

\subsection{Knowledge Graphs for LLMs}\label{sec:rw-kg}

Pan et al.~\cite{pan2024unifying} provide a comprehensive roadmap for
unifying LLMs and knowledge graphs, identifying four paradigms: KG-enhanced
LLMs, LLM-augmented KGs, synergized systems, and KG-only inference.
Structured KG embeddings such as TransE~\cite{bordes2013transe} established
the representation-learning approach to relational data, while KGQA
work~\cite{saxena2020kgqa} demonstrated that KG embeddings can support
complex multi-hop question answering.

More recent approaches~\cite{luo2024rag} use LLMs to reason over KG
triples interactively, traversing the graph via prompted chain-of-thought.
Zhu et al.~\cite{zhu2023large} survey how LLMs are being used as IR systems.
CKG differs from these approaches in representation: rather than full
ontologies or relational triples, it uses lightweight 4-column DAGs trading
expressiveness for construction simplicity and token efficiency.

\subsection{Evaluation Gaps in IR}\label{sec:rw-eval}

Standard IR metrics (F1, MRR, NDCG) do not account for token cost.
RAGAS~\cite{es2023ragas} measures faithfulness and relevance but not
efficiency. Token cost has begun receiving attention as LLM inference
scales: retrieval systems that return large, redundant contexts inflate
cost without improving accuracy. This paper introduces Reasoning Density
Score (RDS $=$ F1 / tokens) and validates it on a 44-domain corpus---the
first metric to jointly optimize quality and token consumption at
multi-domain scale.

\subsection{Educational Knowledge Graphs}\label{sec:rw-edu}

McCreary~\cite{mccreary2024textbooks} introduced the Intelligent Textbooks
methodology, which uses a task-specific data structure we call a
\emph{learning graph} to drive chapter planning, concept sequencing, and
content generation. This paper is the first to formalize this structure as
a citable multi-domain benchmark. Because the learning graph serves a
purpose distinct from that of a general-purpose knowledge graph, and
because the distinction is load-bearing for interpreting our results, we
define the term precisely here and contrast it with its namesake.

\paragraph{Learning graph, defined.}
A \emph{learning graph} (LG) is a small, domain-scoped DAG whose nodes are
atomic teachable concepts and whose edges encode a recommended learning
order (prerequisite relationships). A lightweight taxonomy assigns each
concept to one of a small number of categories (typically 10--16, the
practical upper bound at which a human viewer can reliably distinguish
node colors in a graph visualization), used primarily to color nodes so
that a subject-matter expert can recognize structural patterns at a
glance. A complete formal
definition is given in Section~\ref{sec:lg-def} as
Definition~\ref{def:learning-graph}; for the purposes of related-work
comparison it suffices to note four characteristics: (i) the node set
represents a single bounded pedagogical domain (e.g., high-school geometry,
introductory calculus), not an open-world entity inventory; (ii) edges
encode one relation type---prerequisite---not an arbitrary schema of
relations; (iii) the graph is deliberately small (200--550 concepts in the
McCreary corpus) so that a single human expert can review and maintain it
in its entirety; and (iv) taxonomy categories are a visualization aid,
not a formal ontology.

\paragraph{Contrast with general-purpose knowledge graphs.}
Table~\ref{tab:lg-vs-kg} summarizes the differences between an LG and a
general-purpose KG of the kind represented by DBpedia, Wikidata, YAGO, or
enterprise ontologies. The structures share a graph-theoretic surface but
diverge on nearly every design axis that matters for downstream use.

\begin{table}[htbp]
\centering
\small
\caption{Learning Graph (LG) vs.\ general-purpose Knowledge Graph (KG).}
\label{tab:lg-vs-kg}
\renewcommand{\arraystretch}{1.15}
\begin{tabular}{p{0.22\textwidth}p{0.34\textwidth}p{0.34\textwidth}}
\toprule
\textbf{Dimension} & \textbf{Learning Graph (LG)} & \textbf{General-purpose KG} \\
\midrule
Primary purpose & Accelerate textbook and course content generation;
                  guide concept sequencing. &
                  Answer open-world queries; power search, chatbots, and
                  question answering. \\
Scope           & Single bounded pedagogical domain. &
                  Open-world or large enterprise domain. \\
Typical size    & 200--550 concepts. &
                  $10^5$ to $10^9$+ entities. \\
Node semantics  & Atomic teachable concept. &
                  Named entity, class, property, or literal. \\
Edge semantics  & One relation: recommended learning order (prerequisite). &
                  Many relation types: \texttt{is-a}, \texttt{part-of},
                  \texttt{located-in}, \texttt{born-in}, and dozens more. \\
Schema / ontology & Lightweight taxonomy of 10--16 categories, primarily
                    for visualization. &
                    Formal ontology (RDFS, OWL, SHACL) with class hierarchy,
                    datatype constraints, and inference rules. \\
Structural constraint & Must be acyclic (valid teaching order must exist). &
                        Cycles are permitted and common. \\
Quality criterion & Pedagogical fidelity: an expert teacher agrees the
                    ordering is sensible. &
                    Factual accuracy: assertions match ground truth about
                    the world. \\
Maintenance model & One SME reviews the entire graph. &
                    Many editors, automated extraction pipelines, versioned
                    releases. \\
\bottomrule
\end{tabular}
\end{table}

\paragraph{Design intent vs.\ benchmark use.}
The McCreary learning graphs were not originally designed as retrieval
structures. They were designed to accelerate textbook content generation:
the DAG dictates chapter order, the taxonomy drives visual navigation,
and the concept labels seed prompts for automated chapter drafting. That
the same structure also serves as a highly token-efficient retrieval
substrate---the finding this paper quantifies---is a secondary consequence
of three properties the original design happened to enforce: a closed
concept vocabulary, deterministic edge traversal, and a corpus-wide
uniform schema. We consider this repurposing itself to be a finding: a
structure built for \emph{content production} turns out to be
competitive as a structure for \emph{structured retrieval}, at four
orders of magnitude lower token cost than dynamically extracted
graph-RAG baselines. The benchmark in this paper evaluates a use case
for which the underlying data structure was not originally intended.

\paragraph{Related educational-graph work.}
Automated prerequisite learning from educational data~\cite{liu2016mooc,
shen2018kdd} has sought to derive concept dependency graphs from MOOCs
and course catalogs---the inverse of CKG's assumption. Where those works
\emph{infer} prerequisite structure from data, CKG \emph{uses}
expert-authored structure directly, trading automation for precision.
Adaptive hypermedia systems~\cite{brusilovsky2001adaptive} have long used
structured concept maps to personalize educational content delivery.
Our benchmark provides the first token-efficiency evaluation of such
structures in the context of LLM-based retrieval.
